\centerline{\bf \Huge Геом диагностика}

\begin{flushright}
        \scriptsize{это должен не уметь каждый десятиклассник}
\end{flushright}

\problem{1}{
В треугольнике $ABC$ угол при вершине $A$ равен $60^\circ$.
Обозначим через $O$, $I$ и $H$ центр описанной окружности,
центр вписанной окружности и ортоцентр треугольника соответственно.
Докажите, что окружность $(HOI)$ проходит через точку $B$.
}

\medskip

\problem{2}{
Дан неравнобедренный треугольник $ABC$, в котором $AB<AC$.
Точка $S$~--- середина дуги $BAC$, точка $P$~--- проекция точки
$S$ на прямую $AC$. Докажите, что точка $P$ делит пополам
периметр ломаной $BAC$.
}

\medskip

\problem{3}{
Высоты остроугольного треугольника $ABC$ пересекаются в точке $H$.
На отрезках $BH$ и $CH$ отмечены точки $B_1$ и $C_1$ соответственно
так, что $B_1C_1\parallel BC$. Оказалось, что центр окружности
$(B_1HC_1)$ лежит на прямой $BC$. Докажите, что окружности
$(ABC)$ и $(B_1HC_1)$ касаются.
}

\medskip

\problem{4}{
Треугольник $ABC$ вписан в окружность $\omega$ и описан вокруг
окружности с центром в точке $I$. Прямые $BI$ и $CI$ пересекают
$\omega$ в точках $X$ и $Y$ соответственно. Прямая $\ell$ проходит
через точку $I$, касается окружности $(XIY)$ и пересекает прямую
$XY$ в точке $S$. Докажите, что прямая $SA$ касается окружности
$(ABC)$.
}

\medskip

\problem{5}{
Дан неравнобедренный треугольник $ABC$, в котором
$\angle B=135^\circ$. Пусть $M$~--- середина отрезка $AC$.
Точка $O$~--- центр окружности $\Omega$, описанной около
треугольника $ABC$. Луч $BM$ вторично пересекает окружность
$\Omega$ в точке $D$. Докажите, что центр окружности $(BOD)$
лежит на прямой $AC$.
}

\medskip

\problem{6}{
Диагонали $AC$ и $BD$ вписанного четырёхугольника $ABCD$
пересекаются в точке $P$. Точка $Q$ выбрана на отрезке $BC$
так, что $PQ\perp AC$. Докажите, что прямая, проходящая через
центры окружностей, описанных около треугольников $APD$ и $BQD$,
параллельна прямой $AD$.
}

\medskip

\problem{7}{
Продолжения боковых сторон $AB$ и $CD$ трапеции $ABCD$
пересекаются в точке $P$. На отрезке $AD$ нашлась такая точка
$Q$, что $BQ=CQ$. Докажите, что линия центров окружностей,
описанных около треугольников $AQC$ и $BQD$, перпендикулярна
прямой $PQ$.
}

\medskip

\problem{8}{
В неравнобедренном треугольнике $ABC$ проведена биссектриса
$BL$. Продолжение медианы, проведённой из вершины $B$,
пересекает окружность $(ABC)$ в точке $D$. Через центр окружности
$(BDL)$ проведена прямая $\ell$, параллельная $AC$. Докажите,
что $\ell$ касается окружности $(ABC)$.
}